% Options for packages loaded elsewhere
% Options for packages loaded elsewhere
\PassOptionsToPackage{unicode}{hyperref}
\PassOptionsToPackage{hyphens}{url}
\PassOptionsToPackage{dvipsnames,svgnames,x11names}{xcolor}
%
\documentclass[
  english,
  11pt,
]{article}
\usepackage{xcolor}
\usepackage[margin=25mm]{geometry}
\usepackage{amsmath,amssymb}
\setcounter{secnumdepth}{-\maxdimen} % remove section numbering
\usepackage{iftex}
\ifPDFTeX
  \usepackage[T1]{fontenc}
  \usepackage[utf8]{inputenc}
  \usepackage{textcomp} % provide euro and other symbols
\else % if luatex or xetex
  \usepackage{unicode-math} % this also loads fontspec
  \defaultfontfeatures{Scale=MatchLowercase}
  \defaultfontfeatures[\rmfamily]{Ligatures=TeX,Scale=1}
\fi
\usepackage{lmodern}
\ifPDFTeX\else
  % xetex/luatex font selection
\fi
% Use upquote if available, for straight quotes in verbatim environments
\IfFileExists{upquote.sty}{\usepackage{upquote}}{}
\IfFileExists{microtype.sty}{% use microtype if available
  \usepackage[]{microtype}
  \UseMicrotypeSet[protrusion]{basicmath} % disable protrusion for tt fonts
}{}
\makeatletter
\@ifundefined{KOMAClassName}{% if non-KOMA class
  \IfFileExists{parskip.sty}{%
    \usepackage{parskip}
  }{% else
    \setlength{\parindent}{0pt}
    \setlength{\parskip}{6pt plus 2pt minus 1pt}}
}{% if KOMA class
  \KOMAoptions{parskip=half}}
\makeatother
% Make \paragraph and \subparagraph free-standing
\makeatletter
\ifx\paragraph\undefined\else
  \let\oldparagraph\paragraph
  \renewcommand{\paragraph}{
    \@ifstar
      \xxxParagraphStar
      \xxxParagraphNoStar
  }
  \newcommand{\xxxParagraphStar}[1]{\oldparagraph*{#1}\mbox{}}
  \newcommand{\xxxParagraphNoStar}[1]{\oldparagraph{#1}\mbox{}}
\fi
\ifx\subparagraph\undefined\else
  \let\oldsubparagraph\subparagraph
  \renewcommand{\subparagraph}{
    \@ifstar
      \xxxSubParagraphStar
      \xxxSubParagraphNoStar
  }
  \newcommand{\xxxSubParagraphStar}[1]{\oldsubparagraph*{#1}\mbox{}}
  \newcommand{\xxxSubParagraphNoStar}[1]{\oldsubparagraph{#1}\mbox{}}
\fi
\makeatother


\usepackage{longtable,booktabs,array}
\newcounter{none} % for unnumbered tables
\usepackage{calc} % for calculating minipage widths
% Correct order of tables after \paragraph or \subparagraph
\usepackage{etoolbox}
\makeatletter
\patchcmd\longtable{\par}{\if@noskipsec\mbox{}\fi\par}{}{}
\makeatother
% Allow footnotes in longtable head/foot
\IfFileExists{footnotehyper.sty}{\usepackage{footnotehyper}}{\usepackage{footnote}}
\makesavenoteenv{longtable}
\usepackage{graphicx}
\makeatletter
\newsavebox\pandoc@box
\newcommand*\pandocbounded[1]{% scales image to fit in text height/width
  \sbox\pandoc@box{#1}%
  \Gscale@div\@tempa{\textheight}{\dimexpr\ht\pandoc@box+\dp\pandoc@box\relax}%
  \Gscale@div\@tempb{\linewidth}{\wd\pandoc@box}%
  \ifdim\@tempb\p@<\@tempa\p@\let\@tempa\@tempb\fi% select the smaller of both
  \ifdim\@tempa\p@<\p@\scalebox{\@tempa}{\usebox\pandoc@box}%
  \else\usebox{\pandoc@box}%
  \fi%
}
% Set default figure placement to htbp
\def\fps@figure{htbp}
\makeatother



\ifLuaTeX
\usepackage[bidi=basic,shorthands=off]{babel}
\else
\usepackage[bidi=default,shorthands=off]{babel}
\fi
\ifLuaTeX
  \usepackage{selnolig} % disable illegal ligatures
\fi


\setlength{\emergencystretch}{3em} % prevent overfull lines

\providecommand{\tightlist}{%
  \setlength{\itemsep}{0pt}\setlength{\parskip}{0pt}}



 


\makeatletter
\@ifpackageloaded{caption}{}{\usepackage{caption}}
\AtBeginDocument{%
\ifdefined\contentsname
  \renewcommand*\contentsname{Table of contents}
\else
  \newcommand\contentsname{Table of contents}
\fi
\ifdefined\listfigurename
  \renewcommand*\listfigurename{List of Figures}
\else
  \newcommand\listfigurename{List of Figures}
\fi
\ifdefined\listtablename
  \renewcommand*\listtablename{List of Tables}
\else
  \newcommand\listtablename{List of Tables}
\fi
\ifdefined\figurename
  \renewcommand*\figurename{Figure}
\else
  \newcommand\figurename{Figure}
\fi
\ifdefined\tablename
  \renewcommand*\tablename{Table}
\else
  \newcommand\tablename{Table}
\fi
}
\@ifpackageloaded{float}{}{\usepackage{float}}
\floatstyle{ruled}
\@ifundefined{c@chapter}{\newfloat{codelisting}{h}{lop}}{\newfloat{codelisting}{h}{lop}[chapter]}
\floatname{codelisting}{Listing}
\newcommand*\listoflistings{\listof{codelisting}{List of Listings}}
\makeatother
\makeatletter
\makeatother
\makeatletter
\@ifpackageloaded{caption}{}{\usepackage{caption}}
\@ifpackageloaded{subcaption}{}{\usepackage{subcaption}}
\makeatother
\usepackage{bookmark}
\IfFileExists{xurl.sty}{\usepackage{xurl}}{} % add URL line breaks if available
\urlstyle{same}
% fallback for those not using the hyperref driver hyperxmp:
\makeatletter
\@ifundefined{xmpquote}{\newcommand{\xmpquote}[1]{#1}}{}
\makeatother
\hypersetup{
  pdftitle={What should happen at the centre of a vortex?},
  pdflang={en},
  colorlinks=true,
  linkcolor={blue},
  filecolor={Maroon},
  citecolor={Blue},
  urlcolor={Blue},
  pdfcreator={LaTeX via pandoc}}


\title{What should happen at the centre of a vortex?}
\usepackage{etoolbox}
\makeatletter
\providecommand{\subtitle}[1]{% add subtitle to \maketitle
  \apptocmd{\@title}{\par {\large #1 \par}}{}{}
}
\makeatother
\subtitle{Building a smooth starting shape before worrying about the
whole flow}
\author{}
\date{2026-09-12}
\begin{document}
\maketitle


\href{../paper-map/article.html}{Start with the paper map} · Main route,
article 2 of 9

A drawing of a whirlpool often shows arrows circling a central point.
But what should an arrow right at that point do? There is no single
``around the centre'' direction there.

That small question reveals a major requirement of the construction. The
fluid must be perfectly smooth at every time before the designated
singular time. We cannot build the final blow-up by putting a hidden
defect at the centre from the start.

This article explains how the paper begins with a regular inner profile.
The key ideas are the circumference of a circle, a smooth graph, and an
approximation that improves as more shape information is included.

\subsection{Use a roundabout as the first
picture}\label{use-a-roundabout-as-the-first-picture}

Imagine dots painted on a rigid rotating disk. Every dot completes one
turn in the same time. A dot farther from the centre travels a longer
circle, so it moves faster.

The circumference is \(2\pi r\), where \(r\) is the distance from the
centre. With a common rotation period, the speed is proportional to
\(r\). A dot halfway out moves half as fast. At the very centre, its
circular speed is zero.

This does not describe the whole vortex in the paper. It gives a simple
model of regular behaviour near a rotation axis: the circulating speed
can vanish smoothly even while the rotation is strong.

\begin{figure}

\centering{

\pandocbounded{\includegraphics[keepaspectratio,alt={Concentric circles accompany a straight line showing circular speed increasing from zero with distance from the centre.}]{figures/reader.png}}

}

\caption{\label{fig-centre}A rigid-disk example has circulating speed
proportional to distance from the centre. The graph reaches zero at the
centre; no circulating direction has to be assigned there. This is an
illustration of regularity, not the paper's computed velocity profile.}

\end{figure}%

Fluid can still move \textbf{along} the axis. Zero circular speed at the
centre does not require the full velocity to be zero.

\subsection{Describe the shape before changing its
size}\label{describe-the-shape-before-changing-its-size}

A \textbf{profile} is simply a graph describing how a quantity varies
with position. Here we need profiles for circular motion, axial motion,
and pressure.

The construction separates the profile from its changing physical
scales. Imagine printing the same picture at different sizes and
changing the colour scale used for speed. The profile describes the
underlying picture; the time-dependent scales describe how narrow and
fast it becomes.

This separation helps organize the equations. A shrinking coordinate
ruler follows the core, so the profile does not have to be recomputed as
an unrelated picture at every time.

The actual radial and axial rulers change at slightly different rates.
For this first explanation, remembering that there are two rulers is
more useful than memorizing their exponents.

\subsection{Let the equations determine the nearby
shape}\label{let-the-equations-determine-the-nearby-shape}

Suppose we know a smooth graph's value at a point. To draw it nearby, we
might first use a constant value, then add a tilt, then add curvature.

An elementary example is

\[
y=1+x+x^2.
\]

At \(x=0.1\), its three terms are \(1\), \(0.1\), and \(0.01\). This
illustrates how powers can organize information close to the starting
point. It is not the vortex formula.

The inner solver uses this general idea. It represents the unknown
radial profiles by powers of the radial coordinate. The fluid equations
relate the coefficients, so the nearby shape is constrained by the
balance of motion, pressure, and viscosity.

Including more coefficients gives a way to check whether the
approximation is settling down. If successive approximations agree
better over a declared region, that is useful numerical evidence. It is
not automatically a guarantee beyond that region.

In the paper, Appendix B supplies analytic existence and continuation
arguments. Our numerical study illustrates how a local coefficient
construction behaves for stated inputs.
\href{https://cdn.openai.com/pdf/32d9f210-8b73-45e0-91bc-82a30aef8a9a/navier-stokes.pdf}{Paper,
Section 4 and Appendix B}

\subsection{Why the upward and downward motions are not perfectly
symmetric}\label{why-the-upward-and-downward-motions-are-not-perfectly-symmetric}

A perfectly symmetric sketch is attractive, but it can make two useful
mechanisms weak in the same place.

The paper needs enough shear near the edge of the inner flow to amplify
the later pulses. Circular shear helps in much of the region; axial
shear supplies another contribution. A small axial bias separates
locations where those mechanisms would otherwise both become
ineffective.

Think of arranging two supports so their weak spots do not line up. The
analogy explains the design choice; the actual condition is a
quantitative statement about the velocity profiles.

This is one reason the flow cannot be specified by ``draw a nice
symmetrical whirlpool.'' Its shape is chosen to support later steps of
the argument.
\href{https://cdn.openai.com/pdf/32d9f210-8b73-45e0-91bc-82a30aef8a9a/navier-stokes.pdf\#page=5}{Paper,
Section 2.1}

\subsection{A local profile still needs a connection to the
outside}\label{a-local-profile-still-needs-a-connection-to-the-outside}

A good centre is one piece of the construction. The profile must extend
far enough, reach a suitable state at its outer edge, and connect to the
annulus and exterior.

Pressure is especially important. We cannot choose its inner value
arbitrarily and assume that the surrounding rotating fluid will fit it.

Our first inner calculation used declared test pressure data to check
the numerical method. That was a useful component check. It was not an
assembled version of the paper's entire vortex.

The next stage supplied pressure from the source's outer schedule. The
distinction is similar to testing a bridge component in a workshop
before calculating the loads it will carry in the actual bridge.

\subsection{Where this fits in the bigger
picture}\label{where-this-fits-in-the-bigger-picture}

\textbf{The smooth inner profile gives the concentrating vortex a
legitimate starting shape before the singular time.} Its edge must also
provide the shear needed by the later wave mechanism.

The growth of the final maximum speed comes from the evolving
concentration, not an initially broken rotation axis. Continue with
\href{../d6_axis_pressure/article.html}{why the centre pressure depends
on the outside}.

The \href{technical-article.html}{technical article},
\href{../../code/d6_inner_profile/study.ipynb}{notebook}, and
\href{source-notes.md}{source notes} give the coupled equations,
numerical convergence checks, and limits of the local calculation.




\end{document}
